\chapter{Fungsi Linear dan Fungsi Afin}\label{ch:g9:linfunc}

``Tiga kilogram berharga tiga kali harga satu kilogram'': itulah
\omterm{def:g9:linfunc:linear}{perbandingan senilai}, cara paling sederhana dua besaran dapat
berhubungan. \omterm{def:g9:linfunc:linear}{Fungsi linear} adalah \omterm{def:g9:linfunc:linear}{perbandingan senilai} dalam bahasa
fungsi; \omterm{def:g9:linfunc:affine}{fungsi afin} menambahkan sejumlah tetap sebagai awalannya.
Grafiknya berupa garis lurus, dan membaca garis itulah keterampilan
yang dibangun bab ini.

\section{Perbandingan senilai dan fungsi linear}

\begin{definition}[Fungsi linear]\label{def:g9:linfunc:linear}
\emph{Fungsi linear}\index{fungsi linear} mengalikan setiap bilangan
dengan sebuah bilangan tetap $a$:
\[
f(x) = a x .
\]
Bilangan $a$ adalah \emph{koefisien} $f$. Dua besaran bersifat
\emph{sebanding}\index{perbandingan senilai} persis ketika yang satu
merupakan fungsi linear dari yang lain.
\end{definition}

\begin{example}\label{ex:g9:linfunc:linear}
Kalau apel berharga $3$ tiap kilogram, harga $x$ kilogram adalah
$f(x) = 3x$: membeli dua kali lipat berharga dua kali lipat, dan harga
$2.5$ kg adalah $f(2.5) = 7.5$.

Sebaliknya, kalau sebuah \omterm{def:g9:linfunc:linear}{fungsi linear} memenuhi $f(4) = 10$, maka
$a = \frac{10}{4} = 2.5$ dan $f(x) = 2.5x$: \emph{satu nilai
menentukan sebuah \omterm{def:g9:linfunc:linear}{fungsi linear} sepenuhnya.}
\end{example}

\begin{proposition}[Grafik fungsi linear]\label{prop:g9:linfunc:lineargraph}
Grafik $f(x) = ax$ adalah garis lurus \emph{yang melalui titik asal}.
Sebaliknya, setiap garis tidak tegak yang melalui titik asal adalah
grafik sebuah \omterm{def:g9:linfunc:linear}{fungsi linear}.
\end{proposition}
\admitted

\begin{omfigure}
\begin{tikzpicture}
\begin{axis}[omaxis,
  width=7.4cm, height=5.6cm,
  xmin=-3.4, xmax=3.4, ymin=-3.2, ymax=3.6,
  xtick={-2,-1,1,2}, ytick={-2,-1,1,2}]
  \addplot[omDef, thick, domain=-2.2:2.2] {1.5*x};
  \addplot[omThm, thick, domain=-3.1:3.1] {-0.5*x};
  \node[omDef, font=\small] at (axis cs:2.15,2.6) {$y = 1.5x$};
  \node[omThm, font=\small] at (axis cs:-2.5,1.7) {$y = -0.5x$};
\end{axis}
\end{tikzpicture}

{\small Dua \omterm{def:g9:linfunc:linear}{fungsi linear}: kedua grafiknya melalui titik asal;
koefisien $a$ adalah tinggi yang tercapai di $x = 1$.}
\end{omfigure}

\section{Fungsi afin}

\begin{definition}[Fungsi afin]\label{def:g9:linfunc:affine}
\emph{Fungsi afin}\index{fungsi afin} berbentuk
\[
f(x) = a x + b ,
\]
dengan $a$ dan $b$ bilangan tetap. \omterm{def:g9:linfunc:linear}{Fungsi linear} adalah keadaan
khususnya ketika $b = 0$; fungsi tetap adalah keadaan $a = 0$.
\end{definition}

\begin{example}
Sebuah sanggar senam menarik biaya pendaftaran $20$, lalu $5$ tiap
kunjungan: biaya seluruhnya untuk $x$ kunjungan adalah
$f(x) = 5x + 20$. Biayanya \emph{tidak} \omterm{def:g9:linfunc:linear}{sebanding} dengan banyaknya
kunjungan (10 kunjungan tidak berharga dua kali lipat harga 5
kunjungan), tetapi tiap kunjungan tambahannya menambah \omterm{def:g1:addition:def}{jumlah} yang
sama, yaitu $5$.
\end{example}

\begin{proposition}[Grafik fungsi afin]\label{prop:g9:linfunc:affinegraph}
Grafik $f(x) = ax + b$ adalah garis lurus:
\begin{itemize}
  \item $b = f(0)$ adalah \emph{titik potong sumbu-$y$}: garisnya
        memotong sumbu tegaknya di $(0, b)$;
  \item $a$ adalah \emph{gradiennya}\index{gradien}: bergerak satu
        satuan ke kanan menggeser $a$ satuan ke atas (ke bawah kalau
        $a < 0$); fungsinya naik ketika $a > 0$, turun ketika $a < 0$;
  \item untuk dua masukan mana pun $u \neq v$,
        $a = \dfrac{f(v) - f(u)}{v - u}$.
\end{itemize}
\end{proposition}

\begin{proof}[Bukti bagian terakhirnya]
$f(v) - f(u) = (av + b) - (au + b) = a(v - u)$; bagilah dengan
$v - u$. Sisanya diterima tanpa bukti pada tingkat ini (buku Sekolah
Menengah memberi pembahasan lengkap tentang garis dan \omterm{def:g8:equations:def}{persamaannya}).
\end{proof}

\begin{omfigure}
\begin{tikzpicture}
\begin{axis}[omaxis,
  width=7.8cm, height=5.8cm,
  xmin=-1.9, xmax=4.4, ymin=-1.4, ymax=5.4,
  xtick={1,2,3}, ytick={1,2,3}]
  \addplot[omDef, thick, domain=-1.6:4.1] {x + 1};
  \fill (axis cs:0,1) circle (1.5pt);
  \node[left, font=\small] at (axis cs:0,1.25) {$b$};
  \draw[dashed, omThm] (axis cs:1,2) -- (axis cs:2,2) -- (axis cs:2,3);
  \node[font=\footnotesize, below] at (axis cs:1.5,2) {$+1$};
  \node[font=\footnotesize, right] at (axis cs:2,2.5) {$+a$};
\end{axis}
\end{tikzpicture}

{\small Membaca $y = ax + b$ pada grafiknya: garisnya memotong
sumbu-$y$ di ketinggian $b$, lalu menanjak $a$ untuk tiap langkah satu
satuan ke kanan (di sini $a = 1$, $b = 1$).}
\end{omfigure}

\begin{method}[Mencari fungsi afin dari dua nilai]\label{met:g9:linfunc:findaffine}
Diberikan $f(u)$ dan $f(v)$ dengan $u \neq v$:
\begin{enumerate}
  \item hitunglah \omterm{prop:g9:linfunc:affinegraph}{gradiennya} $a = \dfrac{f(v) - f(u)}{v - u}$;
  \item sulihkan salah satu dari kedua nilai yang diketahui ke dalam
        $f(x) = ax + b$ untuk mencari $b$;
  \item periksalah dengan nilai yang satunya.
\end{enumerate}
\end{method}

\begin{example}\label{ex:g9:linfunc:findaffine}
Carilah \omterm{def:g9:linfunc:affine}{fungsi afin} dengan $f(2) = 7$ dan $f(5) = 16$.
\[
a = \frac{16 - 7}{5 - 2} = \frac93 = 3 ,
\]
lalu $f(2) = 3 \times 2 + b = 7$ memberi $b = 1$: jadi
$f(x) = 3x + 1$. Periksa: $f(5) = 15 + 1 = 16$.
\end{example}

\section{Persen}

\begin{proposition}[Perubahan persen]\label{prop:g9:linfunc:percent}
Menaikkan sebuah besaran sebesar $t\,\%$ mengalikannya dengan
$1 + \dfrac{t}{100}$; menurunkannya sebesar $t\,\%$ mengalikannya
dengan $1 - \dfrac{t}{100}$. Perubahan persen adalah \omterm{def:g9:linfunc:linear}{fungsi linear}.
\end{proposition}

\begin{proof}
Menaikkan $x$ sebesar $t\,\%$ berarti menambahkan $\frac{t}{100}\,x$:
\[
x + \frac{t}{100}\,x = \left(1 + \frac{t}{100}\right) x .
\]
Perhitungan yang sama dengan tanda minus untuk penurunan.
\end{proof}

\begin{example}\label{ex:g9:linfunc:percent}
Harga $80$ naik $15\%$: harga barunya
$80 \times 1.15 = 92$. Harga $60$ turun $30\%$:
$60 \times 0.7 = 42$.

\emph{Merangkai perubahan berarti mengalikan faktornya.} Kenaikan
$20\%$ yang disusul penurunan $20\%$ memberi
$1.2 \times 0.8 = 0.96$: yaitu \emph{penurunan} $4\%$ secara
keseluruhan, bukan kembali ke awalnya!
\end{example}

\section{Latihan}

\begin{exercise}[$\star$]\label{exo:g9:linfunc:1}
Di antara fungsi ini, mana yang linear? mana yang afin? mana yang bukan keduanya?
\[
f(x) = 4x, \qquad
g(x) = 3x - 5, \qquad
h(x) = x^2, \qquad
k(x) = -\tfrac x2, \qquad
m(x) = 7 .
\]
\end{exercise}

\begin{exercise}[$\star$]\label{exo:g9:linfunc:2}
Misalkan $f(x) = 2x - 3$. Hitunglah $f(0)$, $f(4)$, $f(-1)$, lalu
carilah bilangan $x$ yang membuat $f(x) = 9$.
\end{exercise}

\begin{exercise}[$\star$]\label{exo:g9:linfunc:3}
Sebuah \omterm{def:g9:linfunc:linear}{fungsi linear} memenuhi $f(6) = 15$. Carilah koefisiennya, lalu
hitunglah $f(10)$ dan $f(-2)$.
\end{exercise}

\begin{exercise}[$\star$]\label{exo:g9:linfunc:4}
Gambarlah pada sistem koordinat yang sama grafik $f(x) = 2x$,
$g(x) = 2x + 3$ dan $h(x) = -x + 4$. Apa yang dapat dikatakan tentang
grafik $f$ dan $g$?
\end{exercise}

\begin{exercise}[$\star$]\label{exo:g9:linfunc:5}
Carilah \omterm{def:g9:linfunc:affine}{fungsi afin} yang memenuhi $f(1) = 5$ dan $f(3) = 11$; lalu
fungsi yang memenuhi $g(0) = 4$ dan $g(2) = 0$.
\end{exercise}

\begin{exercise}[$\star\star$]\label{exo:g9:linfunc:6}
Hitunglah: harga sesudah kenaikan $25\%$ atas $120$; harga sesudah
potongan $40\%$ atas $65$; harga aslinya kalau sebuah barang berharga
$69$ sesudah potongan $25\%$.
\end{exercise}

\begin{exercise}[$\star\star$]\label{exo:g9:linfunc:7}
Paket telepon A berharga $10$ tiap bulan ditambah $0.05$ tiap menit
percakapan; paket B berharga $25$ tiap bulan dengan percakapan tanpa
batas.
\begin{enumerate}
  \item Nyatakan biaya bulanan paket A sebagai fungsi banyaknya menit
        $x$.
  \item Mulai dari berapa menit tiap bulannya paket B lebih murah?
\end{enumerate}
\end{exercise}

\begin{exercise}[$\star\star$]\label{exo:g9:linfunc:8}
Populasi $2000$ bakteri bertambah $10\%$ tiap jam.
\begin{enumerate}
  \item Berapa bakterinya sesudah satu jam? Sesudah dua jam?
  \item Jelaskan mengapa jawabannya sesudah dua jam bukan $2400$.
\end{enumerate}
\end{exercise}

\begin{exercise}[$\star\star$]\label{exo:g9:linfunc:9}
Grafik sebuah \omterm{def:g9:linfunc:affine}{fungsi afin} melalui titik $(2, 3)$ dan
$(6, 1)$.
\begin{enumerate}
  \item Hitunglah \omterm{prop:g9:linfunc:affinegraph}{gradiennya}. Apakah fungsinya naik atau turun?
  \item Berikan bentuknya, lalu hitunglah masukan yang keluarannya
        $0$.
\end{enumerate}
\end{exercise}

\begin{exercise}[$\star\star\star$]\label{exo:g9:linfunc:10}
Sebuah toko mula-mula menaikkan sebuah harga sebesar $t\,\%$, lalu
menurunkan harga barunya sebesar $t\,\%$.
\begin{enumerate}
  \item Tunjukkan bahwa harga akhirnya sama dengan harga aslinya
        dikalikan $1 - \dfrac{t^2}{10000}$.
  \item Simpulkan bahwa harga akhirnya selalu \emph{lebih rendah}
        daripada harga aslinya (untuk $0 < t \leq 100$), lalu carilah
        perubahan persen keseluruhannya untuk $t = 10$.
\end{enumerate}
\end{exercise}

\section{Soal: Dua termometer}

\begin{problem}[{Soal akhir pekan --- Celsius, Fahrenheit dan fungsi
afin dalam kehidupan sehari-hari: satu suhu terbaca sama pada kedua
skalanya, dan gradien yang tetap adalah tanda tangan kelurusan}]
\label{pb:g9:linfunc:1}
Orang Amerika mendengar ``$68$ derajat'' lalu berpakaian tipis; orang
Eropa mendengarnya lalu menyambar mantel. Kedua termometer dunia
dihubungkan sebuah \omterm{def:g9:linfunc:affine}{fungsi afin} --- dan membangunnya dari dua kenyataan
persis keterampilan pada \cref{met:g9:linfunc:findaffine}. Soal ini
melukis pengubahannya, mencari suhu yang mengerikan tempat kedua
skalanya sepakat, lalu membaca taksi, jangkrik dan lilin dengan
sepasang kacamata yang sama, dan berakhir dengan ujian yang mengenali
kelurusan itu sendiri.

\medskip
\noindent\textbf{Bagian I --- Membangun pengubahannya.}
Dua kenyataan memancang skala Fahrenheit: air membeku pada
$32\,^\circ$F dan mendidih pada $212\,^\circ$F (yaitu pada $0\,^\circ$C
dan $100\,^\circ$C).
\begin{enumerate}
  \item Carilah \omterm{def:g9:linfunc:affine}{fungsi afin} $F = f(C)$ yang mengubah Celsius
        menjadi Fahrenheit (\cref{met:g9:linfunc:findaffine} dengan
        nilai $f(0) = 32$ dan $f(100) = 212$).
  \item Ubahlah menjadi Fahrenheit: hari yang sejuk $20\,^\circ$C;
        suhu tubuh $37\,^\circ$C; dinginnya $-10\,^\circ$C.
  \item Selesaikan untuk $C$ guna membangun pengubahan baliknya, lalu
        pakailah pada $50\,^\circ$F dan pada $98.6\,^\circ$F.
  \item Apakah $f$ fungsi \emph{linear}
        (\cref{def:g9:linfunc:linear})? Ujilah tanda-tandanya:
        berapa $f(0)$, dan apakah melipatduakan bacaan Celsiusnya
        melipatduakan bacaan Fahrenheitnya? Simpulkan dengan kata
        yang tepat (\cref{def:g9:linfunc:affine}).
  \item Teka-teki klasiknya: adakah suhu yang terbaca
        \emph{bilangan yang sama} pada kedua skalanya? Selesaikan
        $f(x) = x$, lalu uraikan arti penyelesaianmu pada
        grafik $f$ (garis yang mana yang dipotong grafiknya di situ?).
\end{enumerate}

\medskip
\noindent\textbf{Bagian II --- \omterm{def:g9:linfunc:affine}{Fungsi afin} di alam liar.}
\begin{enumerate}[resume]
  \item Taksi A menarik $3$ euro ditambah $1.50$ tiap km; taksi B
        menarik $2$ euro tiap km secara rata. Tulislah kedua fungsi
        harganya lalu bandingkan keduanya untuk perjalanan $4$ km dan
        perjalanan $8$ km.
  \item Carilah jarak titik impas yang membuat kedua taksinya
        berbiaya sama, lalu uraikan keadaannya pada sebuah grafik
        (dua garis --- apa yang terjadi pada titik potongnya, sebelum
        titik itu, dan sesudahnya?).
  \item Patokan kasar para naturalis: seekor jangkrik mengerik
        kira-kira $N = 7T - 30$ kali tiap menit pada suhu
        $T\,^\circ$C. Berapa kerikan pada malam
        $20\,^\circ$C? Kamu membilang $82$ kerikan dalam semenit:
        suhu berapa yang diumumkan jangkriknya?
  \item Sebuah telepon yang dibeli $600$ euro kehilangan nilai $15$
        euro tiap bulan: tulislah fungsi nilainya $v(t)$, carilah
        kapan teleponnya bernilai $240$ euro, dan kapan modelnya
        mengatakan teleponnya tidak bernilai apa-apa. Apakah rumusnya
        masih berarti sesudah tanggal itu?
  \item Perpindahan persen adalah \omterm{def:g9:linfunc:linear}{fungsi linear}: $+25\,\%$ adalah
        perkalian dengan $1.25$, dan $-20\,\%$ dengan $0.8$
        (\cref{prop:g9:linfunc:percent}). Susunlah keduanya: apa yang
        dilakukan kenaikan $+25\,\%$ yang disusul potongan
        $-20\,\%$ terhadap sebuah harga? Jelaskan keajaiban kecilnya,
        lalu bandingkan dengan \cref{exo:g9:linfunc:10}.
\end{enumerate}

\medskip
\noindent\textbf{Bagian III --- \omterm{prop:g9:linfunc:affinegraph}{Gradiennya} adalah segalanya.}
\begin{enumerate}[resume]
  \item Sebuah lilin berukuran $20$ cm ketika dinyalakan dan $17$ cm
        setengah jam kemudian. Dengan mengandaikan pelelehan yang
        afin, carilah $h(t)$ (tinggi dalam cm, $t$ dalam menit), lalu
        ramalkan kapan lilinnya mati.
  \item Dua pesepeda menempuh jalan yang sama: jarak yang pertama
        adalah $d_1(t) = 24t$ (km sesudah $t$ jam), yang kedua
        $d_2(t) = 10 + 18t$. Siapa yang memulai dengan keunggulan
        awal, siapa yang mengayuh lebih cepat, dan pada waktu serta
        kilometer berapa yang pertama menyusul yang kedua?
  \item Misalkan $f(x) = 2x + 3$ dan $g(x) = -0.5x + 7$. Hitunglah
        $f(x) + g(x)$ lalu periksalah bahwa hasilnya afin lagi. Berapa
        \omterm{prop:g9:linfunc:affinegraph}{gradien} dan titik potong sebuah \omterm{def:g1:addition:def}{jumlah}, secara umum?
  \item Bahaslah selengkapnya \omterm{def:g8:equations:def}{persamaan} $mx + p = 0$: berapa
        penyelesaiannya ketika $m \neq 0$ (dan yang mana)? Ketika
        $m = 0$ dan $p \neq 0$? Ketika $m = 0$ dan $p = 0$? (Tiap
        keadaannya punya kisah grafiknya: di mana garis bergradien
        $m$ memotong sumbu mendatarnya?)
  \item Ujian kesegarisannya: apakah titik $(1, 5)$, $(3, 9)$,
        $(7, 17)$ terletak pada satu garis lurus? Bagaimana dengan
        $(1, 5)$, $(3, 9)$, $(6, 16)$? Hitunglah laju perubahannya
        antara tiap pasangnya lalu putuskan. Nyatakan pesan moralnya:
        \emph{laju perubahan yang tetap} adalah tanda tangan \omterm{def:g9:linfunc:affine}{fungsi
        afin} --- grafik yang melengkung mengubah lajunya, dan
        mengukur \emph{itu} adalah kisah untuk buku Sekolah Menengah.

\end{enumerate}
\end{problem}
